
% \section{What is TDOA}

\begin{frame}{Time Difference of Arrival Geometry}
	\begin{figure}
	\centering
	\begin{tikzpicture}[scale=0.5,>=Latex]
		
		%----------------------------------------
		% Receiver locations
		%----------------------------------------
		\coordinate (R1) at (2.5,0);
		\coordinate (R2) at (-2.5,0);
		
		%----------------------------------------
		% Source
		%----------------------------------------
		\coordinate (S) at (3.6,3.0);
		
		%----------------------------------------
		% Compute distances
		%----------------------------------------
		\path let \p1 = ($(R1)-(S)$) in
		\pgfextra{\xdef\radius{veclen(\x1,\y1)}};
		
		\path let \p1 = ($(R2)-(S)$) in
		\pgfextra{\xdef\distSRtwo{veclen(\x1,\y1)}};
		
		%----------------------------------------
		% Hyperbola locus
		% Shifted to align with receiver geometry
		%----------------------------------------
		% Receiver locations
		\coordinate (R1) at (2.5,0);
		\coordinate (R2) at (-2.5,0);
		
		% Source location
		\coordinate (S) at (3.6,3.0);
		
		
		%----------------------------------------
		% Compute hyperbola coefficients
		%----------------------------------------
		
		% c = half receiver spacing
		\pgfmathsetmacro{\cval}{2.5}
		
		% Distances from source to receivers
		\pgfmathsetmacro{\distOne}{sqrt((3.6-2.5)^2+(3.0)^2)}
		\pgfmathsetmacro{\distTwo}{sqrt((3.6+2.5)^2+(3.0)^2)}
		
		\pgfmathsetmacro{\aval}{abs(\distTwo-\distOne)/2}
		
		% b = sqrt(c^2-a^2)
		\pgfmathsetmacro{\bval}{sqrt(\cval*\cval-\aval*\aval)}
		
		% midpoint of receivers
		\pgfmathsetmacro{\xcenter}{(2.5+(-2.5))/2}
		
		% Right branch of hyperbola
		\draw[gray,thin,samples=200,smooth,domain=-2:2]
		plot ({\xcenter+\aval*cosh(\x)}, {\bval*sinh(\x)});
		
		%\node[blue] at (1.0,2.2) {Hyperbola};
		
		%----------------------------------------
		% Circle centered at source through R1
		%----------------------------------------
		\draw[gray,thick] (S) circle (\radius);
		
		%----------------------------------------
		% Point where S-R2 line intersects circle
		%----------------------------------------
		\pgfmathsetmacro{\ratio}{\radius/\distSRtwo}
		\coordinate (P) at ($(S)!\ratio!(R2)$);
		
		%----------------------------------------
		% Radius d to nearer receiver
		%----------------------------------------
		\draw[very thick]
		(S)--(R1)
		node[midway,above right] {$d$};
		
		%----------------------------------------
		% Same distance toward R2
		%----------------------------------------
		\draw[very thick]
		(S)--(P);
		
		%----------------------------------------
		% Excess path length
		%----------------------------------------
		\draw[very thick,dashed,red]
		(P)--(R2)
		node[midway,below] {$\Delta d$};
		
		%----------------------------------------
		% Points
		%----------------------------------------
		\fill (R1) circle (2pt);
		\fill (R2) circle (2pt);
		\fill[red] (S) circle (2.5pt);
		
		\node[below] at (R1) {$R_1$};
		\node[below] at (R2) {$R_2$};
		\node[above left] at (S) {S};
		
	\end{tikzpicture}
	\caption{Example TDOA Locus. $S$ is a source, $R$ is a receiver pair, and $\Delta d$ is the difference in arrival times expressed as an equivalent distance. The half hyperbola shows the set of possible source locations for a given $\Delta d$. }
	\end{figure}
\end{frame}

\begin{frame}{Where We Are Going}
	\begin{figure}
		\centering
		\includegraphics[width=1.0 \textwidth]{../../python_trilateration_sandbox/tdoa/loci_from_truth_points_2d.png}
		\caption{Pairwise receiver-emitter TDOA locus (2D)}
	\end{figure}
\end{frame}
